\subsubsection{Descripción de la problemática u oportunidad}\label{sec:problematica}
Los incendios forestales de los Cerros Orientales de Bogotá plantean un problema de interpretación del territorio durante la respuesta, debido a que el avance del fuego depende de condiciones de viento, relieve y vegetación que pueden variar mientras se organiza la atención. En las zonas próximas a la ciudad, esa evolución puede comprometer viviendas, vías y otros elementos expuestos, de modo que conocer la posición observada del incendio resulta necesario, pero no basta para explorar hacia dónde podría propagarse. Las entrevistas realizadas por el equipo identifican necesidades relacionadas con la dirección y velocidad del avance, los accesos, las posibles barreras y la exposición de personas e infraestructura, además de dificultades para recibir y verificar información oportuna. Estas observaciones se utilizan como evidencia cualitativa de los participantes y no como una afirmación de que todas las instituciones carecen de herramientas de apoyo. En particular, el comandante o coordinador de la operación necesita interpretar información cambiante para orientar su juicio, por lo cual la oportunidad del proyecto consiste en acercar resultados de simulación a una consulta comprensible dentro de ese contexto de decisión~\cite{entrevistas2026}.

La dificultad también tiene una dimensión técnica, ya que representar físicamente la interacción entre fuego, combustible y atmósfera demanda datos y capacidad de cálculo que no pueden suponerse disponibles durante un incidente. Los antecedentes sobre los Cerros Orientales muestran, además, que la caracterización territorial y la calibración de combustibles son problemas propios del contexto local, por lo que disponer de un simulador no resuelve automáticamente la preparación de sus entradas~\cite{verano2013,mendez2018}. De esta manera, existe una oportunidad de Ingeniería de Sistemas en la organización de datos, resultados y mecanismos de consulta que permitan aprovechar preparación computacional previa, sin trasladar al usuario operativo la configuración detallada del modelo ni exigirle interpretar directamente sus archivos científicos. La pertinencia de esa integración fue reforzada en la retroalimentación académica del proyecto, que solicita mantener la aplicación como producto central y hacer explícitas sus condiciones de viabilidad~\cite{pafp2026}.
